\section{Álgebra de conjuntos}

\subsection{Complementos, Diferencias y Diferencias Sim\'etricas}


En esta sección, consideraremos conjuntos que sean subconjuntos de un conjunto universo fijo $\uset.$



El \emph{complemento} $A^{C}$ de un conjunto $A$ es el conjunto de elementos que no pertenecen a $A$, es decir 
$$A^{C}=\set{x\in \uset \mid x \notin A}.$$



Algunos textos denotan $A^{C}$ tambi\'en como $A'$ o $\bar{A}.$ 



El \emph{complemento relativo} de un conjunto $B$ con respecto a un conjunto $A$ se define como 
$$
A\minus B = \set{x \mid x \in A, x \notin B}.
$$
De manera equivalente,	$A \minus B = A \cap B^{C}.$

El conjunto $A\minus B$ se lee \texttt{$A$ menos $B$.} Algunos textos lo denotan tambi\'en como $A-B.$  

La \emph{diferencia sim\'etrica} de los conjuntos $A$ y $B$ se define como $$A\symdif B=\left( A\cup B \right)\minus \left( A \cap B \right).$$



\begin{figure}
	\centering
	\includegraphics[width=10cm,keepaspectratio=true]{./md/venn_complemento.png}
	% venn_complemento.png: 0x0 pixel, 300dpi, 0.00x0.00 cm, bb=
	\caption{Complementos, diferencia y diferencia simétrica.}
	\label{fig:0104}
\end{figure}






\subsection{Conjuntos fundamentales}

Supongamos que 
$$
S=A\cup B, \; A\cap B=\emptyset.
$$  Diremos que $S$ es la unión disjunta de $A$ y $B$ y se denota por $$S=A \sqcup B.$$ 




En general $S$ es una unión disjunta de $P_{1}, P_{2},...,P_{n}$ si 
\begin{itemize}
	\item $\displaystyle S=P_{1}\cup P_{2}\cup...\cup P_{n}$ y
	\item $P_{i}\cap P_{j}=\emptyset$ siempre y cuando $i\neq j.$
\end{itemize}


En este caso, escribimos
$$
S=P_{1}\sqcup P_{2}\sqcup...\sqcup P_{n}.
$$



Diremos que $P_{1}, P_{2},...,P_{n}$ es sistema de conjuntos fundamentales para $\uset$ si
$$
\uset = P_{1}\sqcup P_{2}\sqcup...\sqcup P_{n}.
$$






\begin{figure}
	\centering
	\includegraphics[width=10cm,keepaspectratio=true]{./md/sistema_fundamental.png}
	% sistema_fundamental.png: 0x0 pixel, 300dpi, 0.00x0.00 cm, bb=
	\label{fig:0105}
\end{figure}




\subsection{Álgebra de conjuntos, dualidad}


Los conjuntos bajo las operaciones de unión, intersección y complemento satisface varias leyes o identidad, que se enuncian en la siguiente tabla, y son similares a las leyes de lógica.



\begin{figure}
	\centering
	\includegraphics[width=11cm,keepaspectratio=true]{./md/leyes_conjuntos.png}
	\caption{Leyes de Conjuntos}
	\label{fig:leyesconjuntos}
\end{figure}




Cada ley de conjuntos se corresponde con una ley de lógica. Por ejemplo, la ley de DeMorgan:
\begin{align*}
	\left(A \cup B\right)^{C} &= \set{x \mid x\notin(A \cup B)}\\
	&= \set{x \mid x\notin A \wed x\notin B}\\
	&=A^{C}\cap B^{C}
\end{align*}



{Dualidad}
El \emph{principio de dualidad} establece que la equivalencia $E^{*}$ obtenida a partir de una ley de conjuntos $E$ reemplazando
\[ \cup, \cap, \uset, \emptyset\] por
\[ \cap, \cup, \emptyset, \uset\]
sigue siendo una ley de conjuntos.


A la proposición $E^{*}$ se le conoce como dual de $E.$

\subsection{Cardinalidad}

La \emph{cardinalidad} de un conjunto finito $A,$ denotada por $\card{A}$ o $|A|,$ es el número de elementos distintos que contiene.

\begin{problema}
	Si $A = \set{a,b,c,d},$ entonces $\card{A} = 4.$
	Si $B = \set{x \in \N \mid x^{2} < 10},$ entonces $B = \set{0,1,2,3}$ y $\card{B} = 4.$
\end{problema}

\begin{teorema}[Principio de Inclusión-Exclusión para dos conjuntos]
	Sean $A$ y $B$ conjuntos finitos.  Entonces
	$$
	\card{A \cup B} = \card{A} + \card{B} - \card{A \cap B}.
	$$
\end{teorema}

\begin{proof}
	Al sumar $\card{A} + \card{B},$ los elementos de $A \cap B$ se cuentan dos veces.  Restamos $\card{A \cap B}$ para corregir esta duplicación.
\end{proof}

\begin{problema}
	En una encuesta a 100 estudiantes, 60 estudian inglés, 45 estudian francés y 20 estudian ambos idiomas.  ?`Cuántos estudiantes estudian al menos un idioma?
	
	$$\card{I \cup F} = \card{I} + \card{F} - \card{I \cap F} = 60 + 45 - 20 = 85.$$
	
	Por tanto, 85 estudiantes estudian al menos un idioma, y 15 no estudian ninguno.
\end{problema}

\begin{teorema}[Principio de Inclusión-Exclusión para tres conjuntos]
	Sean $A,$ $B$ y $C$ conjuntos finitos.  Entonces
	$$
	\card{A \cup B \cup C} = \card{A} + \card{B} + \card{C} - \card{A \cap B} - \card{A \cap C} - \card{B \cap C} + \card{A \cap B \cap C}.
	$$
\end{teorema}

\begin{corolario}
	Si $A$ y $B$ son disjuntos $(A \cap B = \emptyset),$ entonces
	$$
	\card{A \cup B} = \card{A} + \card{B}.
	$$
	Más generalmente, si $P_{1}, P_{2}, \dots, P_{n}$ forman una partición de $S,$ entonces
	$$
	\card{S} = \sum_{i=1}^{n} \card{P_{i}}.
	$$
\end{corolario}

\begin{problema}
	En un grupo de 50 personas, 30 prefieren café, 25 prefieren té y 15 prefieren jugo.  Además, 10 prefieren café y té, 8 prefieren café y jugo, 5 prefieren té y jugo, y 3 prefieren los tres.  ?`Cuántas personas no prefieren ninguna de estas bebidas?
	
	\begin{align*}
		\card{C \cup T \cup J} &= 30 + 25 + 15 - 10 - 8 - 5 + 3 \\
		&= 50.
	\end{align*}
	
	Todas las 50 personas prefieren al menos una bebida.
\end{problema}

